\documentclass[a4paper,14pt]{extreport} % формат документа

\usepackage{amsmath}
\usepackage{cmap} % поиск в ПДФ
\usepackage[T2A]{fontenc} % кодировка
\usepackage[utf8]{inputenc} % кодировка исходного текста
\usepackage[english,russian]{babel} % локализация и переносы
\usepackage[left = 2cm, right = 1cm, top = 2cm, bottom = 2 cm]{geometry} % поля
\usepackage{listings}
\usepackage{graphicx} % для вставки рисунков
\usepackage{amsmath}
\usepackage{float}
\usepackage{multirow}
\graphicspath{{pictures/}}
\DeclareGraphicsExtensions{.pdf,.png,.jpg}
\newcommand{\anonsection}[1]{\section*{#1}\addcontentsline{toc}{section}{#1}}

\lstset{ %
	language=Lisp,                % Язык программирования 
	numbers=left,                   % С какой стороны нумеровать          
	frame=single,                    % Добавить рамку
}

\begin{document}
\begin{titlepage}

    \begin{table}[H]
        \centering
        \footnotesize
        \begin{tabular}{cc}
            \multirow{8}{*}{\includegraphics[scale=0.35]{bmstu.jpg}}
            & \\
            & \\
            & \textbf{Министерство науки и высшего образования Российской Федерации} \\
            & \textbf{Федеральное государственное бюджетное образовательное учреждение} \\
            & \textbf{высшего образования} \\
            & \textbf{<<Московский государственный технический} \\
            & \textbf{университет имени Н.Э. Баумана>>} \\
            & \textbf{(МГТУ им. Н.Э. Баумана)} \\
        \end{tabular}
    \end{table}

    \vspace{-2.5cm}

    \begin{flushleft}
        \rule[-1cm]{\textwidth}{3pt}
        \rule{\textwidth}{1pt}
    \end{flushleft}

    \begin{flushleft}
        \small
        ФАКУЛЬТЕТ
        \underline{<<Информатика и системы управления>>\ \ \ \ \ \ \ 
        \ \ \ \ \ \ \ \ \ \ \ \ \ \ \ \ \ \ \ \ \ \ \ \ \ \ \ \ \ \ \ 
    \ \ \ \ \ \ \ \ \ \ \ \ \ \ \ } \\
        КАФЕДРА
        \underline{<<Программное обеспечение ЭВМ и
        информационные технологии>>
        \ \ \ \ \ \ \ \ \ \ \ \ \ \ \ \ \ \ \ \ }
    \end{flushleft}

    \vspace{2cm}

    \begin{center}
        \textbf{Лабораторная работа № 7} \\
        \vspace{0.5cm}
        \textbf{} \\
    \end{center}

    \vspace{4cm}

    \begin{flushleft}
        \begin{tabular}{ll}
            \textbf{Дисциплина} & Функциональное и логическое программирование \\
            \textbf{Студент} & Сиденко А.Г. \\
            \textbf{Группа} & ИУ7-63Б \\
            \textbf{Преподаватель} & Толпинская Н.Б.  \\
        \end{tabular}
    \end{flushleft}

    \vspace{4cm}

   \begin{center}
        Москва, 2020 г.
    \end{center}

\end{titlepage}

\begin{enumerate}
\item \textbf{Написать предикат, который принимает два числа-аргумента и возвращает Т, если первое число не меньше второго.}

\begin{lstlisting}
(defun eqb (a b)
  (if (> b a) Nil T) )
\end{lstlisting}

\item \textbf{Какой из следующих двух вариантов предиката ошибочен и почему?}

\begin{lstlisting}
(defun pred1 (x)
  (and (numberp x) (plusp x)))

(defun pred2 (x)
  (and (plusp x) (numberp x)))
\end{lstlisting}

Ошибочен второй, так как перед выполнением операции, не делается проверка на число. 

\item \textbf{Переписать функцию how-alike, приведенную в лекции и не пользующую COND, используя конструкция IF, AND/OR.}

\begin{lstlisting}
(defun how_alike (x y)
  (cond ((or (= x y) (equal x y)) 'the_same)
        ((and (oddp x) (oddp y)) 'both_odd)
        ((and (evenp x) (evenp y)) 'both_even)
        (T 'difference)) )

(defun how_alike_if (x y)
  (if (or (= x y) (equal x y)) 'the_same
        (if (and (oddp x) (oddp y)) 'both_odd
        (if (and (evenp x) (evenp y)) 'both_even 
        'difference))) )
\end{lstlisting}

Результат:

(how$\_$alike 11 30) $\to$ DIFFERENCE

(how$\_$alike$\_$if 11 30) $\to$ DIFFERENCE

(how$\_$alike 3 3) $\to$ THE$\_$SAME

(how$\_$alike$\_$if 3 3) $\to$ THE$\_$SAME

(how$\_$alike 1 3) $\to$ BOTH$\_$ODD

(how$\_$alike$\_$if 1 3) $\to$ BOTH$\_$ODD

(how$\_$alike -2 30) $\to$ BOTH$\_$EVEN

(how$\_$alike$\_$if -2 30) $\to$ BOTH$\_$EVEN

\end{enumerate}
\end{document}